\subsection{Causal Attention Knockout}

Our probing mechanism is the attention knockout of \citet{geva2023dissecting}, in the
per-layer cross-modal form introduced for multimodal LLMs by
\citet{zhang2025crossmodal}; we restate it here for completeness and to fix notation.

Let $\mathbf{x} = (x_1, \dots, x_T)$ be the full token sequence (prompt + generated
answer) and let $p_\theta(x_t \mid x_{<t})$ be the model's next-token probability
under full attention. For a generated answer of length $L$ starting at position $s$,
we define the \textbf{sequence-level probability} as the geometric mean:
\[
  P_\text{full} = \exp\!\left(\frac{1}{L}\sum_{t=s}^{s+L-1} \log p_\theta(x_t \mid x_{<t})\right).
\]

For a given \emph{block type} $b$ (e.g., image$\to$answer) and a layer window
$[l, l+K)$, we construct a modified attention mask $M^{(b,l)}$ that sets the
corresponding query-key attention scores to $-\infty$ within those layers, zeroing
out those attention weights in the softmax. The blocked probability $P_b^{(l)}$ is
computed via a single forward pass with $M^{(b,l)}$.

The \textbf{per-layer information flow} is:
\[
  \text{IF}(l; b) = P_\text{full} - P_b^{(l)}.
\]
A large positive value means the pathway is causal at that layer. We set $K=1$
(single-layer windows) so every decoder layer is probed independently.

\subsection{Block Types}

We study two pathways:
\begin{itemize}
  \item \textbf{\texttt{image\_to\_last}}: blocks attention from all image tokens to
        the answer tokens, measuring how much visual information drives the answer.
  \item \textbf{\texttt{question\_to\_last}}: blocks attention from all question
        tokens to the answer tokens, measuring question-text influence.
\end{itemize}
The \texttt{full\_attention} control (no blocking) always yields $\text{IF}=0$ by
construction, validating the implementation.

\subsection{Derived Metrics}

\paragraph{Image entry layer.}
\[
  l^* = \min\left\{l : \text{IF}(l;\,\texttt{image\_to\_last}) \geq 0.5\cdot
  \max_{l'} \text{IF}(l';\,\texttt{image\_to\_last})\right\}.
\]
This is the first decoder layer where image information contributes at least half its
peak influence. Lower $l^*$ means the model starts integrating visual evidence earlier.

\paragraph{Question dominance ratio (QDR).}
\[
  \text{QDR} = \frac{\bar{\text{IF}}^+(\texttt{question\_to\_last})}
                    {\bar{\text{IF}}^+(\texttt{image\_to\_last}) + \epsilon},
\]
where $\bar{\text{IF}}^+ = \frac{1}{N}\sum_l \max(\text{IF}(l;\cdot), 0)$.
QDR $> 1$ indicates the model leans more on question text; QDR close to 1 indicates
balanced reliance on both modalities.

\paragraph{Note on absolute vs.\ relative flow.} IF is defined as an \emph{absolute}
drop in sequence probability ($P_\text{full} - P_b^{(l)}$), so a model with higher
baseline confidence can show larger IF for the same proportional reliance. The QDR is a
ratio of two such quantities \emph{within the same model}, which cancels that model's
baseline and makes the across-model QDR comparison robust to it; the IF AUC and
entry-layer metrics are not normalised this way, and we report them as absolute flow.
We also checked the normalised variant (IF divided by $P_\text{full}$, pooled $n{=}250$, descriptive only):
MedGemma's mean normalised flow is about 0.6 times Gemma's in both pathways (image
0.0011 vs.\ 0.0018; question 0.0079 vs.\ 0.0131).

\subsection{Implementation}

The probe is implemented as standard PyTorch \texttt{forward\_pre\_hook}s on each
decoder layer's self-attention module. The hook intercepts the attention mask argument
and applies the block pattern. This makes the framework architecture-agnostic: any
HuggingFace \texttt{AutoModelForImageTextToText} model is supported without
modification. The implementation and the numeric per-layer information-flow exports are
released with the paper (see Section~\ref{sec:conclusion} for the repository).
